\documentclass[12pt,a4paper]{article}
\usepackage[utf8]{inputenc}
\usepackage[T1]{fontenc}
\usepackage{geometry}
\usepackage{xcolor}
\usepackage{titlesec}
\usepackage{enumitem}
\usepackage{amssymb}
\usepackage{helvet}
\usepackage{multicol}
\usepackage{booktabs}
\usepackage{tcolorbox}

\geometry{margin=0.8in}
\renewcommand{\familydefault}{\sfdefault}

% Colors
\definecolor{secblue}{RGB}{20, 60, 120}
\definecolor{seclightblue}{RGB}{240, 245, 255}
\definecolor{secdark}{RGB}{33, 33, 33}
\definecolor{secgrey}{RGB}{120, 120, 120}
\definecolor{dangerred}{RGB}{180, 0, 0}

% Styling
\titleformat{\section}{\color{secblue}\normalfont\huge\bfseries}{\thesection}{1em}{}[\titlerule]
\titleformat{\subsection}{\color{secblue}\normalfont\Large\bfseries}{\thesubsection}{1em}{}

\begin{document}

% --- TITLE PAGE ---
\begin{titlepage}
    \centering
    \vspace*{2cm}
    {\Huge \textbf{sec-npm: The Zero-Trust Sentinel}} \\
    \vspace{0.5cm}
    {\large \textit{A Deep-Dive into Advanced Supply-Chain Defense Architecture}} \\
    \vspace{1.5cm}
    \rule{\linewidth}{2pt} \\
    \vspace{1cm}
    \begin{flushleft}
        \textbf{Project Name:} sec-npm \\
        \textbf{Core Concept:} Passive Observation vs. Active Deception \\
        \textbf{Date:} \today \\
        \textbf{Classification:} Technical Whitepaper
    \end{flushleft}
    \vfill
    \begin{tcolorbox}[colback=seclightblue,colframe=secblue,title=Executive Summary]
        Traditional security tools rely on reactive databases of \textit{known} vulnerabilities. \textbf{sec-npm} introduces a proactive, heuristic-driven engine that identifies malicious intent in real-time. By combining AST de-obfuscation with a high-fidelity shadow sandbox and "Honey-Trap" deception, sec-npm provides a definitive shield against Zero-Day supply-chain hijacks.
    \end{tcolorbox}
\end{titlepage}

\tableofcontents
\newpage

\section{Introduction: The Modern Threat Landscape}
The JavaScript ecosystem is powered by millions of dependencies, creating a massive, interconnected trust network. However, this network is increasingly fragile. The incidents involving \texttt{axios} metadata hijacks and \texttt{ua-parser-js} malware proved that even the most reputable packages can be weaponized in minutes. 

Hackers no longer attack the front door; they attack the \textbf{Supply Chain}. They hijack maintainer accounts, publish malicious updates, and vanish before the community notices. \textbf{sec-npm} was designed to be the "Security Sentry" that interrogates every package before it is allowed to enter the developer's environment.

\section{The Four Pillars of Defense}

\subsection{1. The Similarity Guard (Typosquatting Detection)}
The first line of defense is the \textbf{Levenshtein Similarity Engine}. It analyzes the package name DNA to prevent "Brand Jacking."
\begin{itemize}[label=\checkmark]
    \item \textbf{Heuristic:} If a package name is $\leq 2$ characters away from a top-tier package (e.g., \texttt{axois} vs \texttt{axios}), it is flagged as a Critical Risk.
    \item \textbf{Efficacy:} This stops the most common form of "Spray-and-Pray" attacks where users make simple typing errors.
\end{itemize}

\subsection{2. Deep Static Intelligence (AST \& De-obfuscation)}
Sophisticated malware hides its code using obfuscation. Traditional text-based scanners fail here. \textbf{sec-npm} parses the code into an \textbf{Abstract Syntax Tree (AST)}.
\begin{itemize}
    \item \textbf{Constant Folding:} The tool "solves" string math (e.g., \texttt{'ch' + 'ild\_p' + 'rocess'}) to reveal the true underlying API calls.
    \item \textbf{Entropy Analysis:} We measure the randomness of strings. High-entropy strings often hide encrypted payloads or base64 malware.
    \item \textbf{Binary Detection:} The scanner flags the use of \texttt{Buffer.from()} and Hex-encoded strings, which are often used to construct malicious executables on the fly.
\end{itemize}

\subsection{3. Time-Travel Differential Analysis}
This is the "Historical Memory" of the tool. It compares the current version of a package with its predecessor.
\begin{itemize}
    \item \textbf{The Hijack Fingerprint:} A sudden addition of installation scripts (\texttt{postinstall}, \texttt{preinstall}) in a package that was historically "clean" is the most accurate indicator of a compromised account.
    \item \textbf{Verification:} By checking the "Delta" of the package footprint, we catch hijacks before the malicious script even runs.
\end{itemize}

\subsection{4. Shadow Execution (Active Deception)}
The "Nuclear Option" of sec-npm. Instead of isolation, we use **Deception**.
\begin{itemize}
    \item \textbf{Honey-Traps:} The sandbox is populated with fake credentials, decoy SSH keys, and mock environment variables. 
    \item \textbf{Side-Effect Trapping:} Using ES6 Proxies, we intercept every system call. If a script tries to read a Honey-Trap file or phone home to an external IP, it is caught instantly.
    \item \textbf{Simulated Reality:} The malware thinks it is succeeding, but it is actually performing in a "Digital Theater" where it has no power.
\end{itemize}

\section{Lifecycle Hook Coverage: The Danger Zones}
Malware hides in the "Lifecycle Hooks" that npm runs automatically. \textbf{sec-npm} monitors the entire spectrum:
\begin{itemize}
    \item \textbf{Installation:} \texttt{preinstall}, \texttt{install}, \texttt{postinstall}.
    \item \textbf{Build/Compiling:} \texttt{prepare}, \texttt{prepublish}, \texttt{prepack}.
    \item \textbf{Destruction (The Parting Gift):} \texttt{uninstall}. Malware often tries to wipe the disk when the user discovers the infection and tries to delete the package. We catch this intent in the sandbox before the deletion happens.
\end{itemize}

\section{Operational Philosophy: Heuristics over Blacklists}
A major goal of \textbf{sec-npm} is to avoid hindering legitimate developers. We use a **Scoring System** rather than a binary block. However, certain signals are so high-confidence that they trigger an \textbf{Instant Block (Score: 100)}:
\begin{itemize}
    \item \textbf{Honey-Trap Violation:} Any attempt to access decoy credential files in the sandbox.
    \item \textbf{Critical Typosquatting:} Direct name-jacking of top-tier modules (e.g., \texttt{axois} vs \texttt{axios}).
    \item \textbf{Intelligent Scoring:} For all other cases, we only block when multiple signals overlap (e.g., \textit{New Script} + \textit{Obfuscation} + \textit{Network Access}).
    \item \textbf{Human in the Loop:} The user receives a clear "Green/Yellow/Red" report, allowing them to make the final call with full technical visibility.
\end{itemize}

\section{Case Study: Auditing a Reputable Package (Axios)}
To demonstrate the tool's precision, we conducted an audit of \textbf{axios@1.16.0}, a widely trusted networking library. The results provide a perfect example of the ``Heuristic Intelligence'' model.

\begin{tcolorbox}[colback=seclightblue,colframe=secblue,title=Audit Result: axios@1.16.0]
    \textbf{Risk Score:} 50 / 100 \quad \textbf{Status:} \color{orange}WARNING
    \begin{itemize}
        \item \textbf{Detection:} Identified \texttt{prepare: husky} lifecycle hook.
        \item \textbf{Findings:} 11 AST anomalies detected (Buffer.from, process.env).
        \item \textbf{Shadow Result:} 0 Dynamic Anomalies.
    \end{itemize}
\end{tcolorbox}

\textbf{Security Interpretation:} 
The tool correctly identified that \texttt{axios} uses sensitive patterns (accessing environment variables for proxy settings and using Buffer for data payloads). While these patterns are common in malware, the **Shadow Execution** confirmed that no malicious intent was present (no Honey-Trap triggers, no unauthorized network calls). 

The score of 50 serves as a ``Cautionary Flag.'' For a networking library, this is expected. However, if these same patterns were found in a simple UI component or a utility library, the tool's warning would be the critical signal that prevents a major breach.

\section{Conclusion}
\textbf{sec-npm} represents a shift from "Detection" to "Deception." By creating a Zero-Trust environment where code must prove its innocence through behavioral simulation, we have built a fortress that is faster, safer, and more intelligent than traditional sandboxing.

\textit{Don't just install. Secure.}

\end{document}
